\documentclass[11pt,a4paper]{article}
\usepackage{fontspec}
\usepackage{geometry}
\usepackage{amsmath,amssymb,mathtools}
\usepackage{booktabs}
\usepackage{enumitem}
\usepackage{xcolor}
\usepackage{titlesec}
\usepackage{fancyhdr}
\usepackage[unicode=true]{hyperref}
\usepackage[most]{tcolorbox}
\usepackage{lastpage}
\usepackage{microtype}
\usepackage{xeCJK}
\geometry{top=22mm,bottom=24mm,left=24mm,right=24mm,headheight=14pt}
\definecolor{ink}{HTML}{172033}\definecolor{accent}{HTML}{315A8A}\definecolor{accentlight}{HTML}{EAF1F8}\definecolor{rulegray}{HTML}{D7DEE8}\definecolor{muted}{HTML}{5C6678}
\setmainfont{DejaVu Serif}\setsansfont{DejaVu Sans}\setmonofont{DejaVu Sans Mono}\setCJKmainfont{Noto Serif CJK KR}\setCJKsansfont{Noto Sans CJK KR}
\hypersetup{colorlinks=true,linkcolor=accent,urlcolor=accent,citecolor=accent,pdftitle={모듈로 3 수열의 반주기 대칭성},pdfsubject={수학적 형식화와 알고리즘 감사},pdfkeywords={점화식, 모듈로3, 반주기성, 푸리에, 동반행렬}}
\pagestyle{fancy}\fancyhf{}\fancyhead[L]{\small\color{muted}모듈로 3 반주기 대칭성}\fancyhead[R]{\small\color{muted}형식화와 감사}\fancyfoot[C]{\small\color{muted}\thepage\,/\,\pageref*{LastPage}}\renewcommand{\headrulewidth}{0.3pt}\renewcommand{\footrulewidth}{0pt}
\titleformat{\section}{\Large\bfseries\color{ink}}{\thesection.}{0.6em}{}[\vspace{0.2em}\color{rulegray}\titlerule]\titleformat{\subsection}{\large\bfseries\color{accent}}{\thesubsection}{0.6em}{}\titlespacing*{\section}{0pt}{2.2ex plus .5ex minus .2ex}{1.3ex}\titlespacing*{\subsection}{0pt}{1.7ex plus .4ex minus .2ex}{0.8ex}
\setlist[itemize]{leftmargin=1.6em,itemsep=0.25em,topsep=0.4em}\setlist[enumerate]{leftmargin=1.8em,itemsep=0.25em,topsep=0.4em}\setlength{\parindent}{0pt}\setlength{\parskip}{0.55em}\allowdisplaybreaks
\newtcolorbox{summarybox}{colback=accentlight,colframe=accent,boxrule=0.6pt,arc=2mm,left=4mm,right=4mm,top=3mm,bottom=3mm,before skip=1em,after skip=1em}\newcommand{\F}{\mathbb F}\newcommand{\Z}{\mathbb Z}\newcommand{\concat}{\mathbin{\Vert}}
\begin{document}
\begin{titlepage}\thispagestyle{empty}\vspace*{22mm}{\color{accent}\rule{\textwidth}{1.2pt}}\\[15mm]{\fontsize{27}{33}\selectfont\bfseries\color{ink}모듈로 3 수열의 반주기 대칭성\par}\vspace{9mm}{\Large\color{accent}수학적 형식화와 알고리즘 감사\par}\vspace{13mm}\begin{summarybox}이 문서는 흔히 혼동되는 세 현상, 즉 순환의 부호 반전, 반주기적 반주기성, 역순을 엄밀히 구분한다. 0이 아닌 원시 단어의 유일한 반대 이동을 증명하고 조합적·스펙트럴·행렬적 특성화를 제시하며, 이를 모듈러 점화식의 완전 탐색 분석기와 연결한다.\end{summarybox}\vfill{\large\color{muted}저장소: \texttt{Modular-recurrence-symetry-analyzer-}\par}{\large\color{muted}통합 PDF 판 - 2026년 8월 5일\par}\vspace{12mm}{\color{accent}\rule{\textwidth}{1.2pt}}\end{titlepage}
\renewcommand{\contentsname}{목차}\tableofcontents\clearpage
\section{목적}
이 문서는 저장소 \href{https://github.com/59200/k-bonacci-pisano-finder}{\texttt{59200/k-bonacci-pisano-finder}} 에서 “3에 대한 상보 주기”, “둘로 나누면 자기상보”라고 비형식적으로 표현된 현상을 정식화한다.

불변적인 틀은 체 \(\F_3=\{0,1,2\}\) 위의 주기 단어이며, 덧셈 대합은 \(\nu(x)=-x\pmod3\), 즉 \(\nu(0)=0\), \(\nu(1)=2\), \(\nu(2)=1\) 이다. 따라서 “3에 대한 보수”는 \textbf{모듈로 3 덧셈 역원} 또는 \textbf{모듈로 3 부호 반전 보수}라고 해야 한다. 0은 0으로 가므로 선택한 정수 대표들의 십진 합이 항상 3인 것은 아니다.
\section{서로 다른 세 성질}
\(w=(w_0,\ldots,w_{T-1})\in\F_3^T\) 를 원시 주기 \(T\) 인 단어라 하자.
\subsection{반대 주기의 쌍}
모든 \(n\) 에 대해 \(v_n=-w_n\pmod3\) 이면 \(w\) 와 \(v\) 는 반대 주기의 쌍이다. 두 주기는 서로 다른 순환에 속할 수 있다.
\subsection{반주기적 반주기성}
\(T=2h\) 이고 모든 \(n\) 에 대해 \(w_{n+h}=-w_n\pmod3\) 이면 \(w\) 를 \textbf{반주기적}이라 한다. 이때 \(w=H\concat(-H)\), \(H=(w_0,\ldots,w_{h-1})\) 이다. 예: \(01120221=0112\concat0221\), \(011022=011\concat022\), \(01220211=0122\concat0211\).
\subsection{후반부의 역순}
연산 \(\mathcal R(a_0,\ldots,a_{h-1})=(a_{h-1},\ldots,a_0)\) 은 부호 반전과 독립적이다. 피보나치 예에서 \(\mathcal R(0221)=1220\). 따라서 \(1220\) 은 \(0112\) 의 \textbf{역순 대척 단어}이지 후반부 자체가 아니다.
\section{유일한 반대 이동 정리}
\begin{summarybox}\textbf{정리.} \(w\) 를 주기 \(T\) 인 0이 아닌 원시 단어라 하자. 모든 \(n\) 에 대해 \(w_{n+r}=-w_n\) 인 이동 \(r\) 이 존재하면 \(T\) 는 짝수이고 \(r\equiv T/2\pmod T\) 이다.\end{summarybox}
\subsection*{증명}
이동을 두 번 적용하면 \(w_{n+2r}=w_n\) 이다. 원시성으로부터 \(T\mid2r\) 이다. 0이 아닌 \(\F_3\) 단어는 \(w=-w\) 를 만족할 수 없으므로 \(r\not\equiv0\pmod T\) 이다. 순환군 \(\Z/T\Z\) 의 유일한 비자명한 2차 원소는 \(T\) 가 짝수일 때만 존재하며 \(T/2\) 이다. \hfill\(\square\)
\subsection*{결과}
따라서 0이 아닌 원시 주기는 \begin{enumerate}\item 부호 반전에 의해 다른 순환으로 보내지거나,\item 부호 반전에 불변이며 이 경우 반드시 반주기적이다.\end{enumerate} 이 이분법은 \(00111201,00222102\) 와 같은 반대 순환 쌍과 \(011022\) 같은 내재적 반주기를 구분한다.
\section{조합적 결과}
\(w=H\concat(-H)\) 이면 1과 2의 개수는 같고 0의 개수는 짝수이며 \(\sum_{n=0}^{T-1}w_n=0\) 이 \(\F_3\) 에서 성립한다. 이는 필요조건이지만 충분조건은 아니다. 중심 승격 \(\widetilde0=0\), \(\widetilde1=1\), \(\widetilde2=-1\) 을 쓰면 생성다항식은 \[W(X)=\sum_{n=0}^{2h-1}\widetilde w_nX^n=(1-X^h)\sum_{n=0}^{h-1}\widetilde w_nX^n\] 으로 분해된다.
\section{푸리에 특성화}
\(\widehat w(k)=\sum_{n=0}^{2h-1}\widetilde w_ne^{-2\pi i kn/(2h)}\) 라 하자. 그러면 \(w_{n+h}=-w_n\) 과 모든 짝수 \(k\) 에 대해 \(\widehat w(k)=0\) 인 것은 동치이다. 순방향에서는 인자 \(1-(-1)^k\) 가 나타나 짝수 \(k\) 에서 사라진다. 역으로 모든 짝수 모드가 0이면 단어는 홀수 모드가 생성하는 부분공간에 있고, \(h\) 이동은 \(-I\) 로 작용한다. 이는 단순한 시각 비교가 아니라 정확한 스펙트럴 판정이다.
\section{모듈로 3 선형 수열}
\(k\) 차 점화식 \[u_{n+k}=c_0u_n+\cdots+c_{k-1}u_{n+k-1}\pmod3\] 을 생각하자. 상태는 \(s_n=(u_n,\ldots,u_{n+k-1})^{\mathsf T}\), 진화는 \(s_{n+1}=Ms_n\) 이며 \(M\) 은 동반행렬이다.
\subsection{순수 주기성}
\(M\) 의 행렬식은 부호를 제외하면 \(c_0\) 이다. \(\F_3\) 위에서 \(M\) 이 가역인 것과 \(c_0\neq0\) 인 것은 동치이다. 이 경우 모든 상태 궤도는 순수 주기적이다. \(c_0=0\) 이면 전주기가 나타날 수 있으므로 순환과 구분해야 한다.
\subsection{순환의 부호 반전}
선형성으로부터 \(M(-s)=-M(s)\) 이다. 따라서 부호 반전은 각 순환을 같은 길이의 순환으로 보낸다. 0이 아닌 순환에는 두 경우만 있다: 서로 다른 두 순환이 교환되거나, 하나의 순환이 불변이며 반드시 반주기적이다. 영 순환은 유일한 퇴화 예외로서 각 점이 고정되고 비자명한 원시 반주기를 정의하지 않는다.
\subsection{하나의 초기값에 대한 판정}
초기값 \(s_0\) 에 대해 \(M^hs_0=-s_0\) 이면 상태의 모든 선형 출력은 \(u_{n+h}=-u_n\) 을 만족한다. \(s_0\neq0\) 이고 순환이 원시적이면 이는 위의 반주기를 준다.
\subsection{균일 판정}
모든 초기값이 같은 반주기 관계 \(h\) 를 만족하는 것과 \(M^h=-I\) 는 동치이다. 또한 최소다항식은 \(\mu_M(X)\mid X^h+1\) 을 \(\F_3[X]\) 에서 만족하며, 따라서 \(M^{2h}=I\) 이다.
\section{피보나치 예}
\(M=\begin{pmatrix}0&1\\1&1\end{pmatrix}\) 를 \(\F_3\) 위에서 계산하면 \(M^4=-I\), \(M^8=I\) 이다. 초기값 \((0,1)\) 은 \(01120221=0112\concat(-0112)\) 를 생성한다. 후반부 역순은 \(\mathcal R(0221)=1220\) 이다. 십진 표기의 우연 \(1220=L_{14}+F_{14}=2F_{15}\) 은 추가적인 부호화 항등식일 뿐 반주기성의 일반적 결과가 아니다.
\section{임의의 모듈로 3 수열로의 일반화}
주어진 주기 단어를 분석하는 데 점화식 가정은 필요 없다. 모든 원시 주기에 대해 회전 순환, 덧셈 역원, 아핀 대칭 \(w_{n+r}=aw_n+b\), 역순 대칭 \(w_{r-n}=aw_n+b\), 반주기성, 짝수·홀수 푸리에 모드를 계산할 수 있다. 행렬 이론은 선형 점화식이 명시된 경우에만 필요하다.
\section{일반 모듈러스로의 확장}
\(\Z/m\Z\) 에서도 자연 대합은 \(x\mapsto-x\) 이고 \(w_{n+h}=-w_n\pmod m\) 정의는 그대로 유효하다. \(m=2\) 에서는 부호 반전이 항등이므로 반주기성이 보통 주기성과 같아진다. \(m>2\) 에서는 구분이 비자명하다. “\(m\) 에 대한 보수”는 정수 대표 선택에 의존하므로 대수적 정의로 부적절하며, 불변 용어는 “모듈로 \(m\) 덧셈 역원”이다.
\section{알고리즘 감사}
\(k\) 차 모듈로 \(m\) 점화식의 주기는 유한 상태공간 \((\Z/m\Z)^k\) 에서 검출해야 하며 임의 길이 접두부에서 반복 블록을 찾으면 안 된다. 스크립트는 밑 \(m\) 압축 상태 부호화, 단일 초기값에 대한 Brent 알고리즘(시간 \(O(\mu+\lambda)\), 메모리 \(O(1)\)), 모든 초기값에 대한 정확한 함수 그래프 분해(시간 \(O(m^k)\)), 계수족 병렬화, 반대 순환과 반주기 순환의 별도 분류, \(M^h=-I\) 행렬 검사를 사용한다. Mathematica나 외부 래퍼는 필요 없다.
\section{재현 가능한 스캔 결과}
열다섯 개의 명명된 계수족을 스캔하면 tritetranacci 족에서 길이 24와 8의 서로 다른 반대 순환 쌍과 반주기 \(011022\), \(01220211\), \(12\) 를 얻는다. 두 번째 스캔은 차수 2부터 6까지의 0이 아닌 이진 계수벡터 \(119\) 개를 다룬다: \(\sum_{k=2}^{6}(2^k-1)=119\). 결과는 전역 판정 \(M^h=-I\) 을 만족하는 족 13개, 비자명한 반주기 순환을 적어도 하나 갖는 족 88개, 서로 다른 반대 순환 쌍을 적어도 하나 갖는 족 96개이다. 저장소의 JSON 및 CSV 파일로 재현할 수 있다.
\section*{결론}\addcontentsline{toc}{section}{결론}
모듈로 3 부호 반전은 주기 단어와 선형 점화식 순환 위의 구조적 대합으로 작용한다. 0이 아닌 원시 단어가 이 대합에 불변일 수 있는 경우는 유일한 반주기 이동뿐이다. 같은 성질은 생성다항식의 인수분해, 짝수 푸리에 모드의 소멸, 또는 선형 경우 행렬 관계 \(M^h=-I\) 로 동등하게 판정된다. 따라서 알고리즘 분석은 순환의 내재적 대칭과 반대 순환 사이의 단순한 대응을 명확히 분리한다.
\end{document}